%
% spec-opencode-lifecycle.tex
%
% Session Lifecycle and claim-forked: A Z Specification
% A state machine over the operations of docs/design-opencode-plugin.md
% (DES-078) that concern how a session roster comes to exist, persists, is
% reclaimed, and how a forked child acquires a delegation by a sequence of
% file writes a process may abandon at any point. Bead ethos-5y4r.
%
% This file is the second, independent machine that
% spec-opencode-integration.tex (the Registry machine) could not host:
% probcli reads one document as one B machine, so a second machine placed
% in that file promoted every one of its schemas to an operation of the
% Registry and destroyed the Registry's certified-complete validation.
% The two documents share the vocabulary of three given sets and ZBOOL and
% nothing else; this file reads and writes no Registry component.
%
% Compile: pdflatex spec-opencode-lifecycle.tex
%
% Must be .tex: this repo's fuzz-check gate globs docs/*.tex for the fuzz
% usepackage, and probcli's Z-to-B translation produces 0 states against
% any other extension.
%
\documentclass[a4paper,10pt,fleqn]{article}
\usepackage[margin=1in]{geometry}
\usepackage{fuzz}
\usepackage{booktabs}
\usepackage{tabularx}
\usepackage[colorlinks=true,linkcolor=blue,citecolor=blue,urlcolor=blue]{hyperref}

\newcommand{\fpath}[1]{\texttt{#1}}
\newcommand{\field}[1]{\texttt{#1}}

\begin{document}

\title{Session Lifecycle and \texttt{claim-forked}: A Z Specification}
\author{Mike S, for Punt Labs}
\date{October 2026}
\hypersetup{
  pdftitle={Session Lifecycle and claim-forked: A Z Specification},
  pdfauthor={Punt Labs},
  pdfsubject={Z Specification},
  pdfcreator={fuzz/probcli}
}
\maketitle
\tableofcontents
\newpage

% ===================================================================
\section{Session Lifecycle and \texttt{claim-forked}}
\label{sec:lifecycle}
% ===================================================================


\texttt{spec-opencode-integration.tex} models how a session is
\emph{correlated} with a delegation. It says nothing about how a session
roster comes to exist, how it persists or is reclaimed, or how a forked
child acquires a delegation by a sequence of file writes that a process
may abandon at any point. Those are the subject of
\fpath{docs/design-opencode-plugin.md} (DES-078): \S2.1.1 (\field{ensure},
the two sidecars, \field{isStale}), \S2.3 and \S2.3.1 (the ancestor walk),
\S2.6 (\field{session delete}), and \S5 points~1 and~2 (\field{Store.Ensure}
and \field{claim-forked}). That prose has been through three rounds of
review (rounds 5, 6 and 7 of PR \#536), and each round found a
counterexample the previous one had not. The organisation's rule for a
defect class that survives two rounds is to stop reviewing and model it;
this document is that model.

The model is a state machine over the operations \emph{as the design
document currently states them}. It was written first with every
amendment switched off, so that ProB could be asked whether the invariants
of \S\ref{sec:lc-invariants} hold of the prose as it stood. They did not.
\S\ref{sec:lc-ledger} records each counterexample, the amendment that
removes it, and the switch that turns that amendment on. The operations
below, with every switch on, are the authority the design document is to
be brought into line with; with a switch off, the corresponding operation
is the prose of rounds 6 and 7.

\subsection{Method, and two facts about the toolchain}
\label{sec:lc-method}

\textbf{One machine per file.} \texttt{probcli} reads a \fuzz{} document
as a single B machine: the state is the schema that \field{Init} primes,
and every other schema is promoted to an operation. The Registry machine
of \texttt{spec-opencode-integration.tex} and the machine of this document
are independent, so they live in separate files; an earlier revision
embedded this one in the Registry file and the promotion of its schemas
raised the Registry's transition count from 3862 to 14692 and left its
validation uncertified. This file loads directly in \texttt{fuzz} and
\texttt{probcli}: \field{Init} is the empty system, and the scenario
initialisations \field{LcInitPair} and \field{LcInitChain} are alternatives
for it (\S\ref{sec:lc-init}), selected by the commands of
\S\ref{sec:lc-verification}, as are the slices.

\textbf{Convention for invariants.} Every invariant that is a state
predicate is checked as the negation of that predicate given to ProB as a
\texttt{-goal}, never as a predicate of a state schema. This is stated once
here and applies to every row of \S\ref{sec:lc-verification}. The reason
is the next paragraph.

\textbf{State predicates are guards.} \texttt{probcli}'s Z-to-B
translation puts only the typing of the state schema into the B
\texttt{INVARIANT}; every predicate in the state schema is conjoined to
the post-state of each operation, as a guard. An operation that would
violate such a predicate is not reported as violating it: it is
silently disabled, and the state is simply not reached. A ``zero
invariant violations'' report from \texttt{/z-spec:test} therefore says
nothing about the predicate part of a state schema (the earlier
Validation figures of \texttt{spec-opencode-integration.tex} are of this
kind; see \S\ref{sec:lc-findings}). The state schema of this section accordingly has
\emph{no predicate}. Each invariant is stated as its own schema
(\field{LcInv}\ldots), which \texttt{probcli} never sees, and is checked by
giving ProB its negation as a \texttt{-goal}: a counterexample is then a
reachable state satisfying the goal, and ``no counterexample, all states
visited'' is exhaustive for the instance.

\textbf{Slicing.} The full machine has twenty-nine operations over
twenty-six variables; ProB explores about two hundred states a second.
Each invariant is therefore checked on a \emph{slice}: the same schemas
with the operations that cannot affect that invariant removed, and a
scenario \field{Init} where a particular lineage is needed. A removed
operation leaves its variables at their initial values. The slices are
named in \S\ref{sec:lc-verification}.

\subsection{Types and constants}

Sessions, missions and delegations are the given sets of
\texttt{spec-opencode-integration.tex} (its \S``Basic Types''), repeated
here so that this file stands alone, together with the boolean it uses for
the amendment switches.

\begin{zed}
[SESSIONID, MISSIONID, DELEGATIONID]
\end{zed}

\begin{zed}
ZBOOL ::= ztrue | zfalse
\end{zed}

The liveness of a roster's numeric participant is
abstracted to three kinds, because \field{isStale} distinguishes no more:
a solo or persona-named roster ages out by its \field{Started} time
(\field{rkAged}); a roster whose primary participant is a numeric id is
stale exactly when that process is dead (\field{rkNumLive},
\field{rkNumDead}). A binding's origin is as in \fpath{active.go}:101--119.
Time is $0 \upto \field{lcMaxTime}$ with a stale-after threshold
\field{lcTtl}; a last-seen \emph{later} than \field{now} models a clock that
has stepped backwards, which is the future-dated case.

\begin{zed}
RKIND ::= rkAged | rkNumLive | rkNumDead
\end{zed}

\begin{zed}
BORIGIN ::= boClaim | boDispatch
\end{zed}

\begin{axdef}
lcMaxTime : \nat \\
lcTtl : \nat
\where
lcMaxTime = 2 \\
lcTtl = 1
\end{axdef}

The nine constants below are the amendments of \S\ref{sec:lc-ledger}.
Setting a constant to \field{zfalse} yields the operation as written in
DES-078 rounds 6 and 7; to \field{ztrue}, the amended operation. Each
amended operation is a conjunct guarded by its constant, so the two
versions differ in nothing else.

\begin{axdef}
fixDeleteOrder, fixDeleteCloses, fixCycleCheck, fixRefuseReconciles, fixInertBinding, fixStaleLive, fixCorruptPurge, fixPostCheck, fixSupersede : ZBOOL
\where
fixDeleteOrder = ztrue \\
fixDeleteCloses = ztrue \\
fixCycleCheck = ztrue \\
fixRefuseReconciles = ztrue \\
fixInertBinding = ztrue \\
fixStaleLive = ztrue \\
fixCorruptPurge = ztrue \\
fixPostCheck = ztrue \\
fixSupersede = ztrue
\end{axdef}

\subsection{State}
\label{sec:lc-state}

The state is the conjunction of eighteen small schemas, so that an
operation can name the components it leaves alone with $\Xi$ rather than
by equations. Session level: \field{roster} is the set of existing roster
files, with \field{rkind}, \field{started} and \field{corrupt} (a roster
whose file no longer parses). \field{parentSide} is the
\fpath{parent-session} sidecar; it is a function from a session to its
recorded parent, and may have an entry for a session with no roster.
\field{seen} is the \fpath{last-seen} sidecar and \field{seenBad} the
sessions whose sidecar does not parse. Claim level: \field{claimed} and
\field{corigin} are the active-mission sidecar, \field{binding} the
delegation-binding sidecar, \field{pend} and \field{pendWit} the
\fpath{claim-forked-pending} record (delegation and witness), and
\field{stage} and its companions are a \field{claim-forked} write phase that
is \emph{in flight}, holding the child's dispatch-pending lock (stage $k$
is the next write to perform; \S\ref{sec:lc-ops}). Delegation level:
\field{used} are the ids ever allocated by \field{NewID}, \field{skel} the
skeletons on disk, \field{dOpen} those whose verdict is still open,
\field{dMis} and \field{dOwn} a skeleton's mission and child, and
\field{missOpen} the missions that are open. The last three schemas are
ghosts: they record history so that an invariant over a transition can be
checked as an invariant over a state.

\begin{schema}{LcTime}
now : 0 \upto lcMaxTime
\end{schema}

\begin{schema}{LcRoster}
roster : \power SESSIONID
\end{schema}

\begin{schema}{LcKind}
rkind : SESSIONID \pfun RKIND
\end{schema}

\begin{schema}{LcStarted}
started : SESSIONID \pfun (0 \upto lcMaxTime)
\end{schema}

\begin{schema}{LcCorrupt}
corrupt : \power SESSIONID
\end{schema}

\begin{schema}{LcLineage}
parentSide : SESSIONID \pfun SESSIONID
\end{schema}

\begin{schema}{LcSeen}
seen : SESSIONID \pfun (0 \upto lcMaxTime) \\
seenBad : \power SESSIONID
\end{schema}

\begin{schema}{LcClaim}
claimed : SESSIONID \pfun MISSIONID \\
corigin : SESSIONID \pfun BORIGIN
\end{schema}

\begin{schema}{LcBinding}
binding : SESSIONID \pfun DELEGATIONID
\end{schema}

\begin{schema}{LcPending}
pend : SESSIONID \pfun DELEGATIONID \\
pendWit : SESSIONID \pfun SESSIONID
\end{schema}

\begin{schema}{LcFlight}
stage : SESSIONID \pfun (0 \upto 6) \\
cfD : SESSIONID \pfun DELEGATIONID \\
cfM : SESSIONID \pfun MISSIONID \\
cfW : SESSIONID \pfun SESSIONID
\end{schema}

\begin{schema}{LcUsed}
used : \power DELEGATIONID
\end{schema}

\begin{schema}{LcSkeleton}
skel : \power DELEGATIONID \\
dMis : DELEGATIONID \pfun MISSIONID \\
dOwn : DELEGATIONID \pfun SESSIONID
\end{schema}

\begin{schema}{LcOpen}
dOpen : \power DELEGATIONID
\end{schema}

\begin{schema}{LcMission}
missOpen : \power MISSIONID
\end{schema}

\begin{schema}{LcSealed}
sealed : SESSIONID \pfun \power SESSIONID
\end{schema}

\begin{schema}{LcPurged}
purgedLive : \power SESSIONID
\end{schema}

\begin{schema}{LcLeak}
leaked : \power SESSIONID
\end{schema}

\begin{zed}
Lifecycle \defs LcTime \land LcRoster \land LcKind \land LcStarted \land LcCorrupt \land LcLineage \land LcSeen \land LcClaim \land LcBinding \land LcPending \land LcFlight \land LcUsed \land LcSkeleton \land LcOpen \land LcMission \land LcSealed \land LcPurged \land LcLeak
\end{zed}



\field{sealed} records, when a roster is written, the parent that was then
recorded for it; \field{LineageStable} says the sidecar never departs from
it while the roster exists. \field{purgedLive} collects every session whose
roster a purge removed while it was demonstrably live. \field{leaked}
collects every child for which a \field{claim-forked} call ended with an
interrupted attempt's skeleton still open and no longer resumable.


\subsection{The invariants}
\label{sec:lc-invariants}

The four invariants of the mission contract, each as a schema over
\field{Lifecycle}, followed by the two checks needed to state the first
precisely and the well-formedness conditions that validate the model
itself. A schema here is a statement, not a guard (\S\ref{sec:lc-method}).

\begin{itemize}
  \item \field{LineageStable}: while a roster exists, the parent that
    \field{Store.ParentOf} reports is the one that was recorded when the
    roster was written.
  \item \field{Acyclic}: the recorded-parent relation has no cycle. This is
    the strong form. It is stated of every sidecar, not only of those
    reachable from some roster, and it makes the visited set of the
    ancestor walk (\S2.3.1) a defence against a corrupt file rather than a
    necessity; the walk keeps its visited set and refuses on a repeat.
  \item \field{NoLiveSessionPurged}: no purge removes a roster whose
    last-seen is within the TTL, or whose numeric participant is live.
    ``Live'' is judged from what ethos can read: the participant kind of an
    unreadable roster is not evidence.
  \item \field{NoOrphanOpenSkeleton}: every open fork-propagated skeleton is
    named by a pending record or by the binding of an existing child. This
    literal form is the first conjunct of the stronger property below.
  \item \field{ReconcileOnExit} (second conjunct of the stronger
    \field{NoOrphanOpenSkeleton}): no \field{claim-forked} call ends, other
    than by ``write failed'', leaving an interrupted attempt's skeleton open
    when it has not adopted it and its binding is not live. The literal form
    holds of a crashed attempt for ever, because the record names the
    skeleton; the harm of an unreconciled attempt is that nothing may ever
    reconcile it (the child may be refused for good), so the postcondition
    is on the exit of the call.
  \item \field{SweepComplete}: an open skeleton whose mission is closed is
    named by a pending record. This is what \field{Store.Close}'s sweep is
    for.
\end{itemize}

\begin{schema}{LcInvLineageStable}
Lifecycle
\where
\forall s : roster @ s \in \dom sealed \land sealed(s) = parentSide \limg \{ s \} \rimg
\end{schema}

\begin{schema}{LcInvAcyclic}
Lifecycle
\where
\forall s : \dom parentSide @ s \notin (parentSide \plus) \limg \{ s \} \rimg
\end{schema}

\begin{schema}{LcInvNoLiveSessionPurged}
Lifecycle
\where
purgedLive = \emptyset
\end{schema}

\begin{schema}{LcInvNoOrphanOpenSkeleton}
Lifecycle
\where
\forall d : dOpen @ (\exists c : \dom pend @ pend(c) = d) \lor (\exists c : roster \cap \dom binding @ binding(c) = d)
\end{schema}

\begin{schema}{LcInvReconcileOnExit}
Lifecycle
\where
leaked = \emptyset
\end{schema}

\begin{schema}{LcInvSweepComplete}
Lifecycle
\where
\forall d : dOpen @ (\exists c : \dom pend @ pend(c) = d) \lor (\exists m : missOpen @ (d \mapsto m) \in dMis)
\end{schema}

\begin{schema}{LcInvWellFormed}
Lifecycle
\where
\dom rkind = roster \\
\dom started = roster \\
corrupt \subseteq roster \\
roster \subseteq \dom sealed \\
seenBad \cap \dom seen = \emptyset \\
\dom corigin = \dom claimed \\
\dom pendWit = \dom pend \\
skel \subseteq used \\
dOpen \subseteq skel \\
\dom dMis = skel \\
\dom dOwn = skel \\
\dom cfM = \dom stage \\
\dom cfW = \dom stage \\
\dom cfD \subseteq \dom stage \\
\forall c : \dom stage | stage(c) \geq 3 @ c \in \dom pend \land pend(c) = cfD(c)
\end{schema}



\subsection{Initialisation}
\label{sec:lc-init}

\field{Init} is the empty system. \field{LcInitPair} is a child with a
parent that holds an operator claim; \field{LcInitChain} adds a grandparent.
Both are scenario initialisations for the slices that need a lineage.

\begin{schema}{Init}
Lifecycle'
\where
now' = 0 \\
roster' = \emptyset \\
rkind' = \emptyset \\
started' = \emptyset \\
corrupt' = \emptyset \\
parentSide' = \emptyset \\
seen' = \emptyset \\
seenBad' = \emptyset \\
claimed' = \emptyset \\
corigin' = \emptyset \\
binding' = \emptyset \\
pend' = \emptyset \\
pendWit' = \emptyset \\
stage' = \emptyset \\
cfD' = \emptyset \\
cfM' = \emptyset \\
cfW' = \emptyset \\
used' = \emptyset \\
skel' = \emptyset \\
dMis' = \emptyset \\
dOwn' = \emptyset \\
dOpen' = \emptyset \\
missOpen' = MISSIONID \\
sealed' = \emptyset \\
purgedLive' = \emptyset \\
leaked' = \emptyset
\end{schema}

\begin{schema}{LcInitChain}
Lifecycle'
\where
\exists a, b, c : SESSIONID; m : MISSIONID | a \neq b \land b \neq c \land a \neq c @ roster' = \{a, b, c\} \land rkind' = \{a \mapsto rkAged, b \mapsto rkAged, c \mapsto rkAged\} \land started' = \{a \mapsto 0, b \mapsto 0, c \mapsto 0\} \land parentSide' = \{b \mapsto a, c \mapsto b\} \land sealed' = \{a \mapsto \emptyset, b \mapsto \{a\}, c \mapsto \{b\}\} \land claimed' = \{a \mapsto m\} \land corigin' = \{a \mapsto boClaim\} \\
now' = 0 \\
corrupt' = \emptyset \\
seen' = \emptyset \\
seenBad' = \emptyset \\
binding' = \emptyset \\
pend' = \emptyset \\
pendWit' = \emptyset \\
stage' = \emptyset \\
cfD' = \emptyset \\
cfM' = \emptyset \\
cfW' = \emptyset \\
used' = \emptyset \\
skel' = \emptyset \\
dMis' = \emptyset \\
dOwn' = \emptyset \\
dOpen' = \emptyset \\
missOpen' = MISSIONID \\
purgedLive' = \emptyset \\
leaked' = \emptyset
\end{schema}

\begin{schema}{LcInitPair}
Lifecycle'
\where
\exists a, b : SESSIONID; m : MISSIONID | a \neq b @ roster' = \{a, b\} \land rkind' = \{a \mapsto rkAged, b \mapsto rkAged\} \land started' = \{a \mapsto 0, b \mapsto 0\} \land parentSide' = \{b \mapsto a\} \land sealed' = \{a \mapsto \emptyset, b \mapsto \{a\}\} \land claimed' = \{a \mapsto m\} \land corigin' = \{a \mapsto boClaim\} \\
now' = 0 \\
corrupt' = \emptyset \\
seen' = \emptyset \\
seenBad' = \emptyset \\
binding' = \emptyset \\
pend' = \emptyset \\
pendWit' = \emptyset \\
stage' = \emptyset \\
cfD' = \emptyset \\
cfM' = \emptyset \\
cfW' = \emptyset \\
used' = \emptyset \\
skel' = \emptyset \\
dMis' = \emptyset \\
dOwn' = \emptyset \\
dOpen' = \emptyset \\
missOpen' = MISSIONID \\
purgedLive' = \emptyset \\
leaked' = \emptyset
\end{schema}



\subsection{Operations}
\label{sec:lc-ops}

\paragraph{Environment.} The clock may be set to any time, so it may step
backwards; a numeric participant dies; a roster or a last-seen file
becomes unparseable.

\begin{schema}{LcClockSet}
\Delta Lifecycle \\
\Xi LcRoster \\
\Xi LcKind \\
\Xi LcStarted \\
\Xi LcCorrupt \\
\Xi LcLineage \\
\Xi LcSeen \\
\Xi LcClaim \\
\Xi LcBinding \\
\Xi LcPending \\
\Xi LcFlight \\
\Xi LcUsed \\
\Xi LcSkeleton \\
\Xi LcOpen \\
\Xi LcMission \\
\Xi LcSealed \\
\Xi LcPurged \\
\Xi LcLeak \\
t? : 0 \upto lcMaxTime
\where
now' = t?
\end{schema}

\begin{schema}{LcPidDies}
\Delta Lifecycle \\
\Xi LcTime \\
\Xi LcRoster \\
\Xi LcStarted \\
\Xi LcCorrupt \\
\Xi LcLineage \\
\Xi LcSeen \\
\Xi LcClaim \\
\Xi LcBinding \\
\Xi LcPending \\
\Xi LcFlight \\
\Xi LcUsed \\
\Xi LcSkeleton \\
\Xi LcOpen \\
\Xi LcMission \\
\Xi LcSealed \\
\Xi LcPurged \\
\Xi LcLeak \\
s? : SESSIONID
\where
s? \in roster \\
(s? \mapsto rkNumLive) \in rkind \\
rkind' = rkind \oplus \{ s? \mapsto rkNumDead \}
\end{schema}

\begin{schema}{LcCorruptRoster}
\Delta Lifecycle \\
\Xi LcTime \\
\Xi LcRoster \\
\Xi LcKind \\
\Xi LcStarted \\
\Xi LcLineage \\
\Xi LcSeen \\
\Xi LcClaim \\
\Xi LcBinding \\
\Xi LcPending \\
\Xi LcFlight \\
\Xi LcUsed \\
\Xi LcSkeleton \\
\Xi LcOpen \\
\Xi LcMission \\
\Xi LcSealed \\
\Xi LcPurged \\
\Xi LcLeak \\
s? : SESSIONID
\where
s? \in roster \\
s? \notin corrupt \\
corrupt' = corrupt \cup \{ s? \}
\end{schema}

\begin{schema}{LcCorruptSeen}
\Delta Lifecycle \\
\Xi LcTime \\
\Xi LcRoster \\
\Xi LcKind \\
\Xi LcStarted \\
\Xi LcCorrupt \\
\Xi LcLineage \\
\Xi LcClaim \\
\Xi LcBinding \\
\Xi LcPending \\
\Xi LcFlight \\
\Xi LcUsed \\
\Xi LcSkeleton \\
\Xi LcOpen \\
\Xi LcMission \\
\Xi LcSealed \\
\Xi LcPurged \\
\Xi LcLeak \\
s? : SESSIONID
\where
s? \in \dom seen \\
seen' = \{ s? \} \ndres seen \\
seenBad' = seenBad \cup \{ s? \}
\end{schema}



\paragraph{\field{ensure}} (\S2.1.1, \S5 point~1). \field{LcEnsureCreate} is a
first registration or a registration over a leftover sidecar: the
lineage sidecar is written first, then the roster, then, if \field{--touch},
the last-seen. The three rules of the design appear as guards in order:
a non-empty parent must have a roster unless \field{--parent-gone}; a
recorded parent must equal the supplied one; and the supplied parent, with
\field{par?} empty or a singleton, is then recorded. With
\field{fixCycleCheck} the new edge is refused if the session is already an
ancestor of the supplied parent. \field{LcEnsureLineageOnly} is the crash
between the sidecar and the roster. \field{LcEnsureRefresh} is the existing
roster with \field{--touch}. \field{LcCreateCli} is \field{session create},
which refuses an id whose sidecar records a parent.

\begin{schema}{LcEnsureCreate}
\Delta Lifecycle \\
\Xi LcTime \\
\Xi LcCorrupt \\
\Xi LcClaim \\
\Xi LcBinding \\
\Xi LcPending \\
\Xi LcFlight \\
\Xi LcUsed \\
\Xi LcSkeleton \\
\Xi LcOpen \\
\Xi LcMission \\
\Xi LcPurged \\
\Xi LcLeak \\
s? : SESSIONID \\
par? : \power SESSIONID \\
gone? : ZBOOL \\
touch? : ZBOOL \\
k? : RKIND
\where
s? \notin roster \\
\# par? \leq 1 \\
gone? = ztrue \implies par? \neq \emptyset \\
parentSide \limg \{s?\} \rimg \neq \emptyset \implies par? = parentSide \limg \{s?\} \rimg \\
parentSide \limg \{s?\} \rimg = \emptyset \implies (par? \subseteq roster \lor gone? = ztrue) \\
(fixCycleCheck = ztrue \land parentSide \limg \{s?\} \rimg = \emptyset) \implies s? \notin (par? \cup ((parentSide \plus) \limg par? \rimg)) \\
roster' = roster \cup \{ s? \} \\
rkind' = rkind \oplus \{ s? \mapsto k? \} \\
started' = started \oplus \{ s? \mapsto now \} \\
parentSide' = parentSide \cup (\{ s? \} \cross par?) \\
touch? = ztrue \implies (seen' = seen \oplus \{ s? \mapsto now \} \land seenBad' = seenBad \setminus \{ s? \}) \\
touch? = zfalse \implies (seen' = seen \land seenBad' = seenBad) \\
sealed' = sealed \oplus \{ s? \mapsto par? \}
\end{schema}

\begin{schema}{LcEnsureLineageOnly}
\Delta Lifecycle \\
\Xi LcTime \\
\Xi LcRoster \\
\Xi LcKind \\
\Xi LcStarted \\
\Xi LcCorrupt \\
\Xi LcSeen \\
\Xi LcClaim \\
\Xi LcBinding \\
\Xi LcPending \\
\Xi LcFlight \\
\Xi LcUsed \\
\Xi LcSkeleton \\
\Xi LcOpen \\
\Xi LcMission \\
\Xi LcSealed \\
\Xi LcPurged \\
\Xi LcLeak \\
s? : SESSIONID \\
par? : \power SESSIONID \\
gone? : ZBOOL
\where
s? \notin roster \\
par? \neq \emptyset \\
parentSide \limg \{s?\} \rimg = \emptyset \\
\# par? \leq 1 \\
gone? = ztrue \implies par? \neq \emptyset \\
parentSide \limg \{s?\} \rimg \neq \emptyset \implies par? = parentSide \limg \{s?\} \rimg \\
parentSide \limg \{s?\} \rimg = \emptyset \implies (par? \subseteq roster \lor gone? = ztrue) \\
(fixCycleCheck = ztrue \land parentSide \limg \{s?\} \rimg = \emptyset) \implies s? \notin (par? \cup ((parentSide \plus) \limg par? \rimg)) \\
parentSide' = parentSide \cup (\{ s? \} \cross par?)
\end{schema}

\begin{schema}{LcEnsureRefresh}
\Delta Lifecycle \\
\Xi LcTime \\
\Xi LcRoster \\
\Xi LcKind \\
\Xi LcStarted \\
\Xi LcCorrupt \\
\Xi LcLineage \\
\Xi LcClaim \\
\Xi LcBinding \\
\Xi LcPending \\
\Xi LcFlight \\
\Xi LcUsed \\
\Xi LcSkeleton \\
\Xi LcOpen \\
\Xi LcMission \\
\Xi LcSealed \\
\Xi LcPurged \\
\Xi LcLeak \\
s? : SESSIONID \\
par? : \power SESSIONID
\where
s? \in roster \\
s? \notin corrupt \\
par? = parentSide \limg \{s?\} \rimg \\
seen' = seen \oplus \{ s? \mapsto now \} \\
seenBad' = seenBad \setminus \{ s? \}
\end{schema}

\begin{schema}{LcCreateCli}
\Delta Lifecycle \\
\Xi LcTime \\
\Xi LcCorrupt \\
\Xi LcLineage \\
\Xi LcSeen \\
\Xi LcClaim \\
\Xi LcBinding \\
\Xi LcPending \\
\Xi LcFlight \\
\Xi LcUsed \\
\Xi LcSkeleton \\
\Xi LcOpen \\
\Xi LcMission \\
\Xi LcPurged \\
\Xi LcLeak \\
s? : SESSIONID \\
k? : RKIND
\where
s? \notin corrupt \\
parentSide \limg \{s?\} \rimg = \emptyset \\
s? \notin \dom stage \\
roster' = roster \cup \{ s? \} \\
rkind' = rkind \oplus \{ s? \mapsto k? \} \\
started' = started \oplus \{ s? \mapsto now \} \\
sealed' = sealed \oplus \{ s? \mapsto \emptyset \}
\end{schema}



\paragraph{\field{deleteFiles}} (\S2.6; \fpath{store.go}:278--301). The call
stops at the first failure, so a failed call has performed a prefix of its
steps; \field{done?} is the length of that prefix. The steps are numbered
in the order the code performs them. As written, the last-seen and
parent-session sidecars precede the roster file, so step~7 is the roster
removal and a failure there leaves a roster whose lineage sidecar has gone.
With \field{fixDeleteOrder} the roster is removed at step~6 and the
lineage sidecar last, so the sidecar can outlive the roster (which the
design already treats as authoritative) but never the reverse. With
\field{fixDeleteCloses} step~1 closes, with verdict \field{aborted}, every
open skeleton named by the session's pending record or binding, before
either is cleared.

\paragraph{Purge} (\fpath{store.go}:385--404). A purge deletes a roster that
does not parse, or whose \field{isStale} holds, through the same
\field{deleteFiles}. When a valid last-seen is present, staleness is decided
by the later of \field{Started} and last-seen alone (round 6). With
\field{fixStaleLive} a live numeric participant also protects the roster;
with \field{fixCorruptPurge} an unreadable roster with a fresh last-seen is
left alone.

\begin{schema}{LcDeleteSession}
\Delta Lifecycle \\
\Xi LcTime \\
\Xi LcFlight \\
\Xi LcUsed \\
\Xi LcSkeleton \\
\Xi LcMission \\
\Xi LcPurged \\
\Xi LcLeak \\
s? : SESSIONID \\
done? : 1 \upto 7
\where
s? \notin \dom stage \\
s? \in roster \cup \dom parentSide \cup \dom seen \cup seenBad \cup \dom claimed \cup \dom binding \cup \dom pend \\
fixDeleteCloses = ztrue \implies dOpen' = dOpen \setminus ((pend \limg \{ s? \} \rimg) \cup (binding \limg \{ s? \} \rimg)) \\
fixDeleteCloses = zfalse \implies dOpen' = dOpen \\
done? \geq 2 \implies (claimed' = \{ s? \} \ndres claimed \land corigin' = \{ s? \} \ndres corigin) \\
done? < 2 \implies (claimed' = claimed \land corigin' = corigin) \\
done? \geq 3 \implies binding' = \{ s? \} \ndres binding \\
done? < 3 \implies binding' = binding \\
done? \geq 4 \implies (pend' = \{ s? \} \ndres pend \land pendWit' = \{ s? \} \ndres pendWit) \\
done? < 4 \implies (pend' = pend \land pendWit' = pendWit) \\
done? \geq 5 \implies (seen' = \{ s? \} \ndres seen \land seenBad' = seenBad \setminus \{ s? \}) \\
done? < 5 \implies (seen' = seen \land seenBad' = seenBad) \\
((fixDeleteOrder = ztrue \land done? = 7) \lor (fixDeleteOrder = zfalse \land done? \geq 6)) \implies parentSide' = \{ s? \} \ndres parentSide \\
\lnot ((fixDeleteOrder = ztrue \land done? = 7) \lor (fixDeleteOrder = zfalse \land done? \geq 6)) \implies parentSide' = parentSide \\
((fixDeleteOrder = ztrue \land done? \geq 6) \lor (fixDeleteOrder = zfalse \land done? = 7)) \implies (roster' = roster \setminus \{ s? \} \land rkind' = \{ s? \} \ndres rkind \land started' = \{ s? \} \ndres started \land corrupt' = corrupt \setminus \{ s? \} \land sealed' = \{ s? \} \ndres sealed) \\
\lnot ((fixDeleteOrder = ztrue \land done? \geq 6) \lor (fixDeleteOrder = zfalse \land done? = 7)) \implies (roster' = roster \land rkind' = rkind \land started' = started \land corrupt' = corrupt \land sealed' = sealed)
\end{schema}

\begin{schema}{LcPurgeSession}
\Delta Lifecycle \\
\Xi LcTime \\
\Xi LcFlight \\
\Xi LcUsed \\
\Xi LcSkeleton \\
\Xi LcMission \\
\Xi LcLeak \\
s? : SESSIONID \\
done? : 1 \upto 7
\where
s? \notin \dom stage \\
s? \in roster \cup \dom parentSide \cup \dom seen \cup seenBad \cup \dom claimed \cup \dom binding \cup \dom pend \\
fixDeleteCloses = ztrue \implies dOpen' = dOpen \setminus ((pend \limg \{ s? \} \rimg) \cup (binding \limg \{ s? \} \rimg)) \\
fixDeleteCloses = zfalse \implies dOpen' = dOpen \\
done? \geq 2 \implies (claimed' = \{ s? \} \ndres claimed \land corigin' = \{ s? \} \ndres corigin) \\
done? < 2 \implies (claimed' = claimed \land corigin' = corigin) \\
done? \geq 3 \implies binding' = \{ s? \} \ndres binding \\
done? < 3 \implies binding' = binding \\
done? \geq 4 \implies (pend' = \{ s? \} \ndres pend \land pendWit' = \{ s? \} \ndres pendWit) \\
done? < 4 \implies (pend' = pend \land pendWit' = pendWit) \\
done? \geq 5 \implies (seen' = \{ s? \} \ndres seen \land seenBad' = seenBad \setminus \{ s? \}) \\
done? < 5 \implies (seen' = seen \land seenBad' = seenBad) \\
((fixDeleteOrder = ztrue \land done? = 7) \lor (fixDeleteOrder = zfalse \land done? \geq 6)) \implies parentSide' = \{ s? \} \ndres parentSide \\
\lnot ((fixDeleteOrder = ztrue \land done? = 7) \lor (fixDeleteOrder = zfalse \land done? \geq 6)) \implies parentSide' = parentSide \\
((fixDeleteOrder = ztrue \land done? \geq 6) \lor (fixDeleteOrder = zfalse \land done? = 7)) \implies (roster' = roster \setminus \{ s? \} \land rkind' = \{ s? \} \ndres rkind \land started' = \{ s? \} \ndres started \land corrupt' = corrupt \setminus \{ s? \} \land sealed' = \{ s? \} \ndres sealed) \\
\lnot ((fixDeleteOrder = ztrue \land done? \geq 6) \lor (fixDeleteOrder = zfalse \land done? = 7)) \implies (roster' = roster \land rkind' = rkind \land started' = started \land corrupt' = corrupt \land sealed' = sealed) \\
s? \in roster \\
(s? \in corrupt \land (fixCorruptPurge = zfalse \lor \lnot (\exists t : 0 \upto lcMaxTime @ (s? \mapsto t) \in seen \land now - t \leq lcTtl))) \lor (s? \notin corrupt \land ((s? \in \dom seen \land now - started(s?) > lcTtl \land (\forall t : 0 \upto lcMaxTime | (s? \mapsto t) \in seen @ now - t > lcTtl) \land (fixStaleLive = zfalse \lor rkind(s?) \neq rkNumLive)) \lor (s? \notin \dom seen \land ((rkind(s?) = rkAged \land now - started(s?) > lcTtl) \lor rkind(s?) = rkNumDead)))) \\
(((\exists t : 0 \upto lcMaxTime @ (s? \mapsto t) \in seen \land now - t \leq lcTtl) \lor (s? \notin corrupt \land rkind(s?) = rkNumLive)) \land ((fixDeleteOrder = ztrue \land done? \geq 6) \lor (fixDeleteOrder = zfalse \land done? = 7))) \implies purgedLive' = purgedLive \cup \{ s? \} \\
\lnot (((\exists t : 0 \upto lcMaxTime @ (s? \mapsto t) \in seen \land now - t \leq lcTtl) \lor (s? \notin corrupt \land rkind(s?) = rkNumLive)) \land ((fixDeleteOrder = ztrue \land done? \geq 6) \lor (fixDeleteOrder = zfalse \land done? = 7))) \implies purgedLive' = purgedLive
\end{schema}


\begin{schema}{LcOperatorClaim}
\Delta Lifecycle \\
\Xi LcTime \\
\Xi LcRoster \\
\Xi LcKind \\
\Xi LcStarted \\
\Xi LcCorrupt \\
\Xi LcLineage \\
\Xi LcSeen \\
\Xi LcBinding \\
\Xi LcPending \\
\Xi LcFlight \\
\Xi LcUsed \\
\Xi LcSkeleton \\
\Xi LcOpen \\
\Xi LcMission \\
\Xi LcSealed \\
\Xi LcPurged \\
\Xi LcLeak \\
s? : SESSIONID \\
m? : MISSIONID \\
o? : BORIGIN
\where
s? \in roster \\
s? \notin \dom stage \\
m? \in missOpen \\
claimed' = claimed \oplus \{ s? \mapsto m? \} \\
corigin' = corigin \oplus \{ s? \mapsto o? \}
\end{schema}

\begin{schema}{LcReleaseClaim}
\Delta Lifecycle \\
\Xi LcTime \\
\Xi LcRoster \\
\Xi LcKind \\
\Xi LcStarted \\
\Xi LcCorrupt \\
\Xi LcLineage \\
\Xi LcSeen \\
\Xi LcBinding \\
\Xi LcPending \\
\Xi LcFlight \\
\Xi LcUsed \\
\Xi LcSkeleton \\
\Xi LcOpen \\
\Xi LcMission \\
\Xi LcSealed \\
\Xi LcPurged \\
\Xi LcLeak \\
s? : SESSIONID
\where
s? \in \dom claimed \\
s? \notin \dom stage \\
claimed' = \{ s? \} \ndres claimed \\
corigin' = \{ s? \} \ndres corigin
\end{schema}



\paragraph{\field{claim-forked}} (\S5 point~2). The read phase takes no lock
and writes nothing; \field{LcCfBegin} is a read phase that passes (the
child's roster loads, the witness holds a claim-origin binding, and the
witness is reached by the walk of \S2.3.1: each hop is read from
\field{parentSide} and the session reached must have a readable roster) and
that goes on to take the child's dispatch-pending lock. \field{LcCfRefuse}
is a read-phase refusal. The write phase reconciles the pending record
first (stage~0), in one of three ways: the child already holds a claim
(\field{LcCfReconcileClaimed}), the record is resumable
(\field{LcCfReconcileReuse}), or it is closed and removed
(\field{LcCfReconcileFresh}). It then allocates an id and re-checks that
the mission is open (stage~2), and writes the record, the skeleton, the
binding and the claim, one stage each (3 to 6), before removing the
record. A process may \field{LcCfCrash} at any stage, leaving the files
written so far; or, if it survives a failed write, \field{LcCfAbort}.

With \field{fixRefuseReconciles} a refusal also reconciles. With
\field{fixInertBinding} the claimed case closes the skeleton unless the
binding is live, that is, unless the child's claim is for the skeleton's
mission. With \field{fixPostCheck} the binding is not written unless the
mission is still open: the check that follows the skeleton write, which
closes the window the unlocked sweep leaves. With \field{fixSupersede} a
reconcile that does not resume the record also closes the skeleton named by
the child's binding when that is a different delegation, because the write
phase is about to overwrite the binding.

\begin{schema}{LcCfRefuse}
\Delta Lifecycle \\
\Xi LcTime \\
\Xi LcRoster \\
\Xi LcKind \\
\Xi LcStarted \\
\Xi LcCorrupt \\
\Xi LcLineage \\
\Xi LcSeen \\
\Xi LcClaim \\
\Xi LcBinding \\
\Xi LcFlight \\
\Xi LcUsed \\
\Xi LcSkeleton \\
\Xi LcMission \\
\Xi LcSealed \\
\Xi LcPurged \\
c? : SESSIONID
\where
c? \notin \dom stage \\
(fixRefuseReconciles = zfalse \land pend' = pend \land pendWit' = pendWit \land dOpen' = dOpen \land (c? \in \dom pend \land pend(c?) \in dOpen \implies leaked' = leaked \cup \{ c? \}) \land (\lnot (c? \in \dom pend \land pend(c?) \in dOpen) \implies leaked' = leaked)) \lor (fixRefuseReconciles = ztrue \land leaked' = leaked \land pend' = \{ c? \} \ndres pend \land pendWit' = \{ c? \} \ndres pendWit \land dOpen' = dOpen \setminus (\{ rd : pend \limg \{ c? \} \rimg | \lnot ((c? \mapsto rd) \in binding \land (\exists m : MISSIONID @ (rd \mapsto m) \in dMis \land (c? \mapsto m) \in claimed)) @ rd \}))
\end{schema}

\begin{schema}{LcCfBegin}
\Delta Lifecycle \\
\Xi LcTime \\
\Xi LcRoster \\
\Xi LcKind \\
\Xi LcStarted \\
\Xi LcCorrupt \\
\Xi LcLineage \\
\Xi LcSeen \\
\Xi LcClaim \\
\Xi LcBinding \\
\Xi LcPending \\
\Xi LcUsed \\
\Xi LcSkeleton \\
\Xi LcOpen \\
\Xi LcMission \\
\Xi LcSealed \\
\Xi LcPurged \\
\Xi LcLeak \\
c? : SESSIONID \\
w? : SESSIONID
\where
c? \notin \dom stage \\
c? \in roster \\
c? \notin corrupt \\
w? \in \dom claimed \\
(w? \mapsto boClaim) \in corigin \\
w? \in (((parentSide \rres (roster \setminus corrupt)) \plus) \limg \{ c? \} \rimg) \\
stage' = stage \oplus \{ c? \mapsto 0 \} \\
cfM' = cfM \oplus \{ c? \mapsto claimed(w?) \} \\
cfW' = cfW \oplus \{ c? \mapsto w? \} \\
cfD' = \{ c? \} \ndres cfD
\end{schema}

\begin{schema}{LcCfReconcileClaimed}
\Delta Lifecycle \\
\Xi LcTime \\
\Xi LcRoster \\
\Xi LcKind \\
\Xi LcStarted \\
\Xi LcCorrupt \\
\Xi LcLineage \\
\Xi LcSeen \\
\Xi LcClaim \\
\Xi LcBinding \\
\Xi LcUsed \\
\Xi LcSkeleton \\
\Xi LcMission \\
\Xi LcSealed \\
\Xi LcPurged \\
c? : SESSIONID
\where
(c? \mapsto 0) \in stage \\
c? \in \dom claimed \\
pend' = \{ c? \} \ndres pend \\
pendWit' = \{ c? \} \ndres pendWit \\
stage' = \{ c? \} \ndres stage \\
cfD' = \{ c? \} \ndres cfD \\
cfM' = \{ c? \} \ndres cfM \\
cfW' = \{ c? \} \ndres cfW \\
(fixInertBinding = ztrue \land dOpen' = dOpen \setminus (\{ rd : pend \limg \{ c? \} \rimg | \lnot ((c? \mapsto rd) \in binding \land (\exists m : MISSIONID @ (rd \mapsto m) \in dMis \land (c? \mapsto m) \in claimed)) @ rd \}) \land leaked' = leaked) \lor (fixInertBinding = zfalse \land dOpen' = dOpen \setminus \{ rd : pend \limg \{ c? \} \rimg | (c? \mapsto rd) \notin binding @ rd \} \land (((c? \in \dom pend \land pend(c?) \in dOpen \land (c? \mapsto pend(c?)) \in binding \land \lnot ((c? \mapsto pend(c?)) \in binding \land (\exists m : MISSIONID @ (pend(c?) \mapsto m) \in dMis \land (c? \mapsto m) \in claimed)))) \implies leaked' = leaked \cup \{ c? \}) \land (\lnot ((c? \in \dom pend \land pend(c?) \in dOpen \land (c? \mapsto pend(c?)) \in binding \land \lnot ((c? \mapsto pend(c?)) \in binding \land (\exists m : MISSIONID @ (pend(c?) \mapsto m) \in dMis \land (c? \mapsto m) \in claimed)))) \implies leaked' = leaked))
\end{schema}

\begin{schema}{LcCfReconcileReuse}
\Delta Lifecycle \\
\Xi LcTime \\
\Xi LcRoster \\
\Xi LcKind \\
\Xi LcStarted \\
\Xi LcCorrupt \\
\Xi LcLineage \\
\Xi LcSeen \\
\Xi LcClaim \\
\Xi LcBinding \\
\Xi LcPending \\
\Xi LcUsed \\
\Xi LcSkeleton \\
\Xi LcMission \\
\Xi LcSealed \\
\Xi LcPurged \\
\Xi LcLeak \\
c? : SESSIONID
\where
(c? \mapsto 0) \in stage \\
c? \notin \dom claimed \\
(c? \in \dom pend \land pend(c?) \in dOpen \land (pend(c?) \mapsto cfM(c?)) \in dMis \land pendWit(c?) = cfW(c?) \land cfM(c?) \in missOpen) \\
(fixSupersede = zfalse \land dOpen' = dOpen) \lor (fixSupersede = ztrue \land dOpen' = dOpen \setminus ((binding \limg \{ c? \} \rimg) \setminus \{ pend(c?) \})) \\
cfD' = cfD \oplus \{ c? \mapsto pend(c?) \} \\
(c? \mapsto pend(c?)) \in binding \implies stage' = stage \oplus \{ c? \mapsto 5 \} \\
(c? \mapsto pend(c?)) \notin binding \implies stage' = stage \oplus \{ c? \mapsto 4 \} \\
cfM' = cfM \\
cfW' = cfW
\end{schema}

\begin{schema}{LcCfReconcileFresh}
\Delta Lifecycle \\
\Xi LcTime \\
\Xi LcRoster \\
\Xi LcKind \\
\Xi LcStarted \\
\Xi LcCorrupt \\
\Xi LcLineage \\
\Xi LcSeen \\
\Xi LcClaim \\
\Xi LcBinding \\
\Xi LcUsed \\
\Xi LcSkeleton \\
\Xi LcMission \\
\Xi LcSealed \\
\Xi LcPurged \\
\Xi LcLeak \\
c? : SESSIONID
\where
(c? \mapsto 0) \in stage \\
c? \notin \dom claimed \\
\lnot (c? \in \dom pend \land pend(c?) \in dOpen \land (pend(c?) \mapsto cfM(c?)) \in dMis \land pendWit(c?) = cfW(c?) \land cfM(c?) \in missOpen) \\
(fixSupersede = zfalse \land dOpen' = dOpen \setminus (pend \limg \{ c? \} \rimg)) \lor (fixSupersede = ztrue \land dOpen' = dOpen \setminus ((pend \limg \{ c? \} \rimg) \cup (binding \limg \{ c? \} \rimg))) \\
pend' = \{ c? \} \ndres pend \\
pendWit' = \{ c? \} \ndres pendWit \\
stage' = stage \oplus \{ c? \mapsto 1 \} \\
cfD' = cfD \\
cfM' = cfM \\
cfW' = cfW
\end{schema}

\begin{schema}{LcCfAlloc}
\Delta Lifecycle \\
\Xi LcTime \\
\Xi LcRoster \\
\Xi LcKind \\
\Xi LcStarted \\
\Xi LcCorrupt \\
\Xi LcLineage \\
\Xi LcSeen \\
\Xi LcClaim \\
\Xi LcBinding \\
\Xi LcPending \\
\Xi LcSkeleton \\
\Xi LcOpen \\
\Xi LcMission \\
\Xi LcSealed \\
\Xi LcPurged \\
\Xi LcLeak \\
c? : SESSIONID \\
d? : DELEGATIONID
\where
(c? \mapsto 1) \in stage \\
d? \notin used \\
cfM(c?) \in missOpen \\
used' = used \cup \{ d? \} \\
stage' = stage \oplus \{ c? \mapsto 2 \} \\
cfD' = cfD \oplus \{ c? \mapsto d? \} \\
cfM' = cfM \\
cfW' = cfW
\end{schema}

\begin{schema}{LcCfExitNoRetry}
\Delta Lifecycle \\
\Xi LcTime \\
\Xi LcRoster \\
\Xi LcKind \\
\Xi LcStarted \\
\Xi LcCorrupt \\
\Xi LcLineage \\
\Xi LcSeen \\
\Xi LcClaim \\
\Xi LcBinding \\
\Xi LcPending \\
\Xi LcUsed \\
\Xi LcSkeleton \\
\Xi LcOpen \\
\Xi LcMission \\
\Xi LcSealed \\
\Xi LcPurged \\
\Xi LcLeak \\
c? : SESSIONID
\where
(c? \mapsto 1) \in stage \\
used = DELEGATIONID \lor cfM(c?) \notin missOpen \\
stage' = \{ c? \} \ndres stage \\
cfD' = \{ c? \} \ndres cfD \\
cfM' = \{ c? \} \ndres cfM \\
cfW' = \{ c? \} \ndres cfW
\end{schema}

\begin{schema}{LcCfWriteRecord}
\Delta Lifecycle \\
\Xi LcTime \\
\Xi LcRoster \\
\Xi LcKind \\
\Xi LcStarted \\
\Xi LcCorrupt \\
\Xi LcLineage \\
\Xi LcSeen \\
\Xi LcClaim \\
\Xi LcBinding \\
\Xi LcUsed \\
\Xi LcSkeleton \\
\Xi LcOpen \\
\Xi LcMission \\
\Xi LcSealed \\
\Xi LcPurged \\
\Xi LcLeak \\
c? : SESSIONID
\where
(c? \mapsto 2) \in stage \\
stage' = stage \oplus \{ c? \mapsto 3 \} \\
pend' = pend \oplus \{ c? \mapsto cfD(c?) \} \\
pendWit' = pendWit \oplus \{ c? \mapsto cfW(c?) \} \\
cfD' = cfD \\
cfM' = cfM \\
cfW' = cfW
\end{schema}

\begin{schema}{LcCfWriteSkeleton}
\Delta Lifecycle \\
\Xi LcTime \\
\Xi LcRoster \\
\Xi LcKind \\
\Xi LcStarted \\
\Xi LcCorrupt \\
\Xi LcLineage \\
\Xi LcSeen \\
\Xi LcClaim \\
\Xi LcBinding \\
\Xi LcPending \\
\Xi LcUsed \\
\Xi LcMission \\
\Xi LcSealed \\
\Xi LcPurged \\
\Xi LcLeak \\
c? : SESSIONID
\where
(c? \mapsto 3) \in stage \\
stage' = stage \oplus \{ c? \mapsto 4 \} \\
skel' = skel \cup \{ cfD(c?) \} \\
dOpen' = dOpen \cup \{ cfD(c?) \} \\
dMis' = dMis \oplus \{ cfD(c?) \mapsto cfM(c?) \} \\
dOwn' = dOwn \oplus \{ cfD(c?) \mapsto c? \} \\
cfD' = cfD \\
cfM' = cfM \\
cfW' = cfW
\end{schema}

\begin{schema}{LcCfWriteBinding}
\Delta Lifecycle \\
\Xi LcTime \\
\Xi LcRoster \\
\Xi LcKind \\
\Xi LcStarted \\
\Xi LcCorrupt \\
\Xi LcLineage \\
\Xi LcSeen \\
\Xi LcClaim \\
\Xi LcPending \\
\Xi LcUsed \\
\Xi LcSkeleton \\
\Xi LcOpen \\
\Xi LcMission \\
\Xi LcSealed \\
\Xi LcPurged \\
\Xi LcLeak \\
c? : SESSIONID
\where
(c? \mapsto 4) \in stage \\
fixPostCheck = ztrue \implies cfM(c?) \in missOpen \\
stage' = stage \oplus \{ c? \mapsto 5 \} \\
binding' = binding \oplus \{ c? \mapsto cfD(c?) \} \\
cfD' = cfD \\
cfM' = cfM \\
cfW' = cfW
\end{schema}

\begin{schema}{LcCfWriteClaim}
\Delta Lifecycle \\
\Xi LcTime \\
\Xi LcRoster \\
\Xi LcKind \\
\Xi LcStarted \\
\Xi LcCorrupt \\
\Xi LcLineage \\
\Xi LcSeen \\
\Xi LcBinding \\
\Xi LcPending \\
\Xi LcUsed \\
\Xi LcSkeleton \\
\Xi LcOpen \\
\Xi LcMission \\
\Xi LcSealed \\
\Xi LcPurged \\
\Xi LcLeak \\
c? : SESSIONID
\where
(c? \mapsto 5) \in stage \\
stage' = stage \oplus \{ c? \mapsto 6 \} \\
claimed' = claimed \oplus \{ c? \mapsto cfM(c?) \} \\
corigin' = corigin \oplus \{ c? \mapsto boClaim \} \\
cfD' = cfD \\
cfM' = cfM \\
cfW' = cfW
\end{schema}

\begin{schema}{LcCfFinish}
\Delta Lifecycle \\
\Xi LcTime \\
\Xi LcRoster \\
\Xi LcKind \\
\Xi LcStarted \\
\Xi LcCorrupt \\
\Xi LcLineage \\
\Xi LcSeen \\
\Xi LcClaim \\
\Xi LcBinding \\
\Xi LcUsed \\
\Xi LcSkeleton \\
\Xi LcOpen \\
\Xi LcMission \\
\Xi LcSealed \\
\Xi LcPurged \\
\Xi LcLeak \\
c? : SESSIONID
\where
(c? \mapsto 6) \in stage \\
pend' = \{ c? \} \ndres pend \\
pendWit' = \{ c? \} \ndres pendWit \\
stage' = \{ c? \} \ndres stage \\
cfD' = \{ c? \} \ndres cfD \\
cfM' = \{ c? \} \ndres cfM \\
cfW' = \{ c? \} \ndres cfW
\end{schema}

\begin{schema}{LcCfCrash}
\Delta Lifecycle \\
\Xi LcTime \\
\Xi LcRoster \\
\Xi LcKind \\
\Xi LcStarted \\
\Xi LcCorrupt \\
\Xi LcLineage \\
\Xi LcSeen \\
\Xi LcClaim \\
\Xi LcBinding \\
\Xi LcPending \\
\Xi LcUsed \\
\Xi LcSkeleton \\
\Xi LcOpen \\
\Xi LcMission \\
\Xi LcSealed \\
\Xi LcPurged \\
\Xi LcLeak \\
c? : SESSIONID
\where
c? \in \dom stage \\
stage' = \{ c? \} \ndres stage \\
cfD' = \{ c? \} \ndres cfD \\
cfM' = \{ c? \} \ndres cfM \\
cfW' = \{ c? \} \ndres cfW
\end{schema}

\begin{schema}{LcCfAbort}
\Delta Lifecycle \\
\Xi LcTime \\
\Xi LcRoster \\
\Xi LcKind \\
\Xi LcStarted \\
\Xi LcCorrupt \\
\Xi LcLineage \\
\Xi LcSeen \\
\Xi LcClaim \\
\Xi LcUsed \\
\Xi LcSkeleton \\
\Xi LcMission \\
\Xi LcSealed \\
\Xi LcPurged \\
\Xi LcLeak \\
c? : SESSIONID
\where
stage(c?) \in 3 \upto 5 \\
dOpen' = dOpen \setminus \{ cfD(c?) \} \\
stage(c?) = 5 \implies binding' = \{ c? \} \ndres binding \\
stage(c?) < 5 \implies binding' = binding \\
pend' = \{ c? \} \ndres pend \\
pendWit' = \{ c? \} \ndres pendWit \\
stage' = \{ c? \} \ndres stage \\
cfD' = \{ c? \} \ndres cfD \\
cfM' = \{ c? \} \ndres cfM \\
cfW' = \{ c? \} \ndres cfW
\end{schema}



\paragraph{Missions.} \field{LcCloseLocked} is \field{Store.Close} with the
exclusive mission lock: it cannot run while a write phase holds the shared
lock (stage $\geq 2$ for that mission). \field{LcCloseUnlocked} is the
fallback of \fpath{store.go}:1107--1114, taken when the lock cannot be
acquired: the sweep runs with no exclusion. \field{LcAbandonMission}
refuses while any skeleton of the mission is open; that guard is
\field{countBlockingDelegations}
(\fpath{store.go}:2004--2010) as an observer, and it is why an unaccounted
open skeleton is a defect and not a curiosity.

\begin{schema}{LcCloseLocked}
\Delta Lifecycle \\
\Xi LcTime \\
\Xi LcRoster \\
\Xi LcKind \\
\Xi LcStarted \\
\Xi LcCorrupt \\
\Xi LcLineage \\
\Xi LcSeen \\
\Xi LcClaim \\
\Xi LcBinding \\
\Xi LcPending \\
\Xi LcFlight \\
\Xi LcUsed \\
\Xi LcSkeleton \\
\Xi LcSealed \\
\Xi LcPurged \\
\Xi LcLeak \\
m? : MISSIONID
\where
m? \in missOpen \\
\lnot (\exists c : \dom stage @ stage(c) \geq 2 \land cfM(c) = m?) \\
missOpen' = missOpen \setminus \{ m? \} \\
dOpen' = \{ d : dOpen | (d \mapsto m?) \notin dMis @ d \}
\end{schema}

\begin{schema}{LcCloseUnlocked}
\Delta Lifecycle \\
\Xi LcTime \\
\Xi LcRoster \\
\Xi LcKind \\
\Xi LcStarted \\
\Xi LcCorrupt \\
\Xi LcLineage \\
\Xi LcSeen \\
\Xi LcClaim \\
\Xi LcBinding \\
\Xi LcPending \\
\Xi LcFlight \\
\Xi LcUsed \\
\Xi LcSkeleton \\
\Xi LcSealed \\
\Xi LcPurged \\
\Xi LcLeak \\
m? : MISSIONID
\where
m? \in missOpen \\
missOpen' = missOpen \setminus \{ m? \} \\
dOpen' = \{ d : dOpen | (d \mapsto m?) \notin dMis @ d \}
\end{schema}

\begin{schema}{LcAbandonMission}
\Delta Lifecycle \\
\Xi LcTime \\
\Xi LcRoster \\
\Xi LcKind \\
\Xi LcStarted \\
\Xi LcCorrupt \\
\Xi LcLineage \\
\Xi LcSeen \\
\Xi LcClaim \\
\Xi LcBinding \\
\Xi LcPending \\
\Xi LcFlight \\
\Xi LcUsed \\
\Xi LcSkeleton \\
\Xi LcOpen \\
\Xi LcSealed \\
\Xi LcPurged \\
\Xi LcLeak \\
m? : MISSIONID
\where
m? \in missOpen \\
\lnot (\exists d : dOpen @ (d \mapsto m?) \in dMis) \\
missOpen' = missOpen \setminus \{ m? \}
\end{schema}



\subsection{Counterexamples and amendments}
\label{sec:lc-ledger}

Every counterexample below was found by ProB against the operations as
DES-078 rounds 6 and 7 state them, that is, with the switches \field{zfalse}.
This section was first committed in that form (revision \texttt{e47050a},
eight switches); the ninth, \field{fixSupersede}, was added when C9 appeared
with the other eight on, and C1 to C8 were run again unchanged on the
final text. Each
trace is the shortest breadth-first path to a state satisfying the negated
invariant, in the notation of the operations; ``$a$, $b$'' are sessions,
$m_1$, $m_2$ missions and $d$ a delegation. The first four are the
counterexamples of the three review rounds; they were reproduced, not
assumed. The last five were not known.

\begin{enumerate}
  \item[\textbf{C1}] \textbf{Reconcile skipped on a read-phase refusal.}
    \field{ReconcileOnExit}; \S5 point~2, the order of the read and write
    phases. $b$ is forked from $a$, which holds a claim. \field{LcCfBegin}$(b,a)$;
    \field{LcCfReconcileFresh}; \field{LcCfAlloc}$(b,d)$; \field{LcCfWriteRecord};
    \field{LcCfWriteSkeleton}; \field{LcCfCrash}; \field{LcCfRefuse}$(b)$.
    The record is reconciled only in the write phase (step~2), after the read
    phase has had the chance to refuse; a refusal returns with $d$ open and
    its record in place, and a child that is refused for good never reaches
    the write phase. \textbf{Amendment} \field{fixRefuseReconciles}: every exit
    other than ``write failed'' takes the child's dispatch-pending lock first
    and reconciles. The prose of round 7, ``every call reconciles first'',
    was true of the write phase only.
  \item[\textbf{C2}] \textbf{\field{deleteFiles} removes the parent sidecar
    before a roster removal that then fails.} \field{LineageStable};
    \S2.1.1 and \S2.6. \field{LcEnsureCreate}$(a)$; \field{LcEnsureCreate}$(b,\{a\})$;
    \field{LcDeleteSession}$(b,6)$: the first six steps succeed, the seventh, the
    roster file, fails. $b$ has a roster, no recorded parent, and
    \field{sealed}$(b)=\{a\}$: a later \field{ensure} for $b$ with any parent
    other than $a$ is a lineage conflict, and with none it succeeds, so the
    roster's lineage has changed. (With \field{--parent-gone} the same
    happens in two steps, with $a$ its own parent.) \textbf{Amendment}
    \field{fixDeleteOrder}: remove the roster at step~6 and the
    \fpath{parent-session} sidecar last. A failure at step~7 leaves a sidecar
    with no roster, which the design already treats as authoritative and
    \field{session delete} already clears.
  \item[\textbf{C3}] \textbf{Re-registration after deletion closes a cycle.}
    \field{Acyclic}; \S2.1.1 rule~1 and \S2.3.1. \field{LcEnsureCreate}$(a)$;
    \field{LcEnsureCreate}$(b,\{a\})$; \field{LcDeleteSession}$(a,7)$;
    \field{LcEnsureCreate}$(a,\{b\})$. $b$ is registered, so rule~1 is
    satisfied, and nothing refuses an edge that completes $a\to b\to a$. The
    design's remark that a cycle can only be written through
    \field{--parent-gone} is false; it is also true without deletion, as
    $a\to b$ with $b$ unregistered and \field{--parent-gone}, then
    $b\to a$. \textbf{Amendment} \field{fixCycleCheck}: when a new edge is
    written, refuse if the session is the supplied parent or an ancestor of
    it by the recorded chain, whether or not the ancestors have rosters.
  \item[\textbf{C4}] \textbf{An operator claim of a different mission leaves
    the skeleton open.} \field{ReconcileOnExit}; \S5 point~2, the round-7
    outcome ``record and binding name the same delegation''. As C1 up to
    \field{LcCfWriteBinding}; \field{LcCfCrash}; the operator claims $m_2$,
    different from the skeleton's mission $m_1$: \field{LcOperatorClaim}$(b,m_2)$;
    a later \field{LcCfBegin}$(b,a)$; \field{LcCfReconcileClaimed}$(b)$. The record
    names the binding's delegation, so the call removes the record only. The
    binding is inert ($m_1\neq m_2$) and the skeleton is open, named by
    nothing that will be read again. The prose accepts this (``stays open only
    until its own mission closes''); \field{countBlockingDelegations} does
    not, since the open skeleton blocks \field{Abandon}. \textbf{Amendment}
    \field{fixInertBinding}: in the claimed case, close the skeleton unless
    the binding is live, that is, unless the child's claim is for the
    skeleton's mission.
  \item[\textbf{C5}] \textbf{\field{deleteFiles} on the child orphans its
    skeleton.} \field{NoOrphanOpenSkeleton} (literal form). As C1 up to
    \field{LcCfWriteSkeleton}; \field{LcCfCrash}; \field{LcDeleteSession}$(b,7)$. The
    pending record, the only name the skeleton had, is cleared (step~4) and the
    roster removed; the skeleton stays open. The design says so, ``a skeleton
    it named stays open only until its mission closes''. For a binding the
    same happens after a completed write. \textbf{Amendment}
    \field{fixDeleteCloses}: step~1 of \field{deleteFiles} closes, with
    verdict \field{aborted}, every open skeleton named by the session's record
    or binding. It runs under the child's dispatch-pending lock and takes
    the mission and delegation locks in the order of \S5 point~2; a failure
    to close is a failure of the call, with the roster kept as the retry token.
  \item[\textbf{C6}] \textbf{The unlocked sweep leaves a skeleton the sweep
    can never reach.} \field{SweepComplete}; \fpath{store.go}:1107--1114.
    \field{LcCfBegin}$(b,a)$; \field{LcCfReconcileFresh}; \field{LcCfAlloc};
    \field{LcCfWriteRecord}; \field{LcCloseUnlocked}$(m_1)$ (the exclusive lock could
    not be taken, so the sweep ran beside the write phase);
    \field{LcCfWriteSkeleton}; \field{LcCfWriteBinding}; \field{LcCfWriteClaim};
    \field{LcCfFinish}. The mission is closed and its skeleton is open and
    bound. The code comment calls this ``a narrow TOCTOU window''; nothing
    ever closes the skeleton. \textbf{Amendment} \field{fixPostCheck}: re-read
    the mission after the skeleton is on disk, before the binding. If the
    mission was closed before the skeleton existed, the check sees it; if
    after, the sweep sees the skeleton. On a closed mission the call aborts
    (\field{LcCfAbort}). This closes the window the locked path does not have.
    It is not enough alone: the aborted write can still be crashed, leaving
    a record, and C5 must be fixed too (\S\ref{sec:lc-verification}).
  \item[\textbf{C7}] \textbf{A live process does not protect a roster that
    has an old last-seen.} \field{NoLiveSessionPurged}; \S2.1.1 table, row~2.
    \field{LcEnsureCreate}$(s,\{p\},\text{touch},\text{numeric live})$;
    \field{LcClockSet}$(2)$; \field{LcPurgeSession}$(s,7)$. The table asserts
    ``stale (idle past the TTL decides, even with a live PID)''. The invariant
    of the mission contract says the opposite, and the row is the
    counterexample. \textbf{Amendment} \field{fixStaleLive}: a live numeric
    participant protects the roster whatever the last-seen says. Last-seen
    still decides for a roster that has no live numeric participant, which is
    every OpenCode roster; and a dead numeric participant with a fresh
    last-seen is still not stale. \textbf{Row~2 of the table reverses.}
  \item[\textbf{C8}] \textbf{A roster that does not parse is purged whatever
    its last-seen says.} \field{NoLiveSessionPurged}; \fpath{store.go}:391--397,
    which deletes on a \field{Load} error before \field{isStale} is consulted.
    \field{LcEnsureCreate}$(s,\ldots,\text{touch})$; \field{LcCorruptRoster}$(s)$;
    \field{LcPurgeSession}$(s,7)$. \textbf{Amendment} \field{fixCorruptPurge}: read
    the last-seen first, as \field{isStale} does; an unreadable roster with a
    valid last-seen within the TTL is left alone. (A Claude Code roster that
    does not parse and has no last-seen is purged as today.)
  \item[\textbf{C9}] \textbf{A second propagation overwrites the binding of
    the first.} \field{NoOrphanOpenSkeleton} (literal form), found only
    after C1 to C8 were amended. A completed propagation: \field{LcCfBegin}$(b,a)$;
    reconcile; \field{LcCfAlloc}$(b,d_1)$ and the four writes; \field{LcCfFinish}. Then
    \field{LcReleaseClaim}$(b)$, the operator releasing the child's claim, and a
    second \field{claim-forked}: \field{LcCfBegin}$(b,a)$; reconcile (the record is
    gone); \field{LcCfAlloc}$(b,d_2)$; the writes up to \field{LcCfWriteBinding}. The
    binding now names $d_2$; $d_1$ is open, named by no record and no
    binding. Nothing in DES-078 prevents this: a released child is a child
    with no claim, and the plugin's memo expires. \textbf{Amendment}
    \field{fixSupersede}: a reconcile that does not resume the record also
    closes the skeleton named by the child's current binding.
\end{enumerate}

Two of the four invariants needed more than one amendment, and one
amendment each is necessary: \field{LineageStable} needs only
\field{fixDeleteOrder}, and \field{Acyclic} needs only
\field{fixCycleCheck}. \field{NoLiveSessionPurged} needs \field{fixStaleLive} and
\field{fixCorruptPurge}, each alone leaving the other's counterexample.
\field{NoOrphanOpenSkeleton} needs \field{fixDeleteCloses} and \field{fixSupersede}
for its literal form, and \field{fixRefuseReconciles} with \field{fixInertBinding} for the exit
postcondition; \field{SweepComplete} needs \field{fixDeleteCloses} with
\field{fixPostCheck}.

\subsection{Verification}
\label{sec:lc-verification}

\paragraph{Extraction.} Given the committed file \texttt{F}, a slice is
produced by the commands below, where \texttt{INIT} is one of
\field{LcInit}, \field{LcInitPair}, \field{LcInitChain}, \texttt{V$n$} is
\texttt{ztrue} for each switch to be on and \texttt{zfalse} otherwise,
and \texttt{DROP} is the list of operation schemas the slice removes.

\begin{verbatim}
awk 'BEGIN{p=1; B="\\" "begin"; E="\\" "end"}
 p==1{print} /[\\]begin[{]document[}]/{p=0}
 /^\[SESSIONID, MISSIONID, DELEGATIONID\]$/ || /^ZBOOL ::= ztrue [|] zfalse$/ \
   {print B "{zed}"; print; print E "{zed}"}
 /^%%LIFECYCLE-BEGIN/{r=1;next} /^%%LIFECYCLE-END/{r=0} r{print}
 END{print E "{document}"}' F > S.tex
for n in fixDeleteOrder fixDeleteCloses fixCycleCheck fixRefuseReconciles \
         fixInertBinding fixStaleLive fixCorruptPurge fixPostCheck fixSupersede; do
  sed -E -i "s/^$n = z(true|false)( [\\\\][\\\\])?\$/$n = V_$n\2/" S.tex   # V_n: ztrue or zfalse
done
sed -i '/^[\\]begin{schema}{LcInv/,/^[\\]end{schema}/d' S.tex
for s in OTHER_INIT_1 OTHER_INIT_2 DROPPED_OP ...; do
  sed -i "/^[\\\\]begin{schema}{$s}\$/,/^[\\\\]end{schema}\$/d" S.tex
done
sed -i 's/^\([\\]begin{schema}\){INIT}$/\1{Init}/' S.tex
probcli S.tex -card SESSIONID s -card MISSIONID m -card DELEGATIONID d \
  -p MAX_INITIALISATIONS 1000 -p MAX_OPERATIONS 1000 \
  -goal "GOAL" -mc N -bf -nodead -check_goal
\end{verbatim}

\noindent The switch sed is written once per switch, with the value that
switch is to have. \texttt{-bf} makes the trace the shortest. \texttt{-nodead}
is needed because a slice with the operations that re-register a session
removed deadlocks when every roster is gone, which is not a violation. The
\texttt{GOAL}s, in B, are the negations of the invariant schemas:

\begin{verbatim}
LineageStable        not(!s.(s:roster => (s:dom(sealed) & sealed(s)=parentSide[{s}])))
Acyclic              not(!s.(s:dom(parentSide) => s /: closure1(parentSide)[{s}]))
NoLiveSessionPurged  purgedLive /= {}
NoOrphanOpenSkeleton #d.(d:dOpen & not((#c.(c:dom(pend) & pend(c)=d)) or
                         (#c.(c:roster & c:dom(binding) & binding(c)=d))))
ReconcileOnExit      leaked /= {}
SweepComplete        #d.(d:dOpen & dMis(d) /: missOpen & not(#c.(c:dom(pend) & pend(c)=d)))
\end{verbatim}

\noindent A goal that reports ``no counterexample'' is also followed by
ProB's own message that the goal was ``unknown'' in the root state, which
has no variables; the result is that of the model check. For two goals the
result was cross-checked as the LTL formula
\texttt{G not(\{GOAL\})}: the amended delete slice gives \emph{TRUE}, the
unamended one a counterexample ending in the same state.

\paragraph{Slices.} Operations are removed by name; a removed operation
leaves its variables at their initial values. \emph{lineage} keeps
\field{LcEnsureCreate}, \field{LcEnsureLineageOnly}, \field{LcCreateCli} and
\field{LcDeleteSession} and removes the six other session-level
operations and every claim and mission operation. \emph{purge} keeps all
session-level operations, including \field{LcPurgeSession}, and removes the
claim-forked, claim and mission operations; the variant marked
\emph{no \field{LcCorruptRoster}} also removes that operation. \emph{claim}
keeps only the \field{LcCf}\ldots operations, \field{LcOperatorClaim} and
\field{LcReleaseClaim}. \emph{delete} is \emph{claim} without
\field{LcOperatorClaim}, \field{LcReleaseClaim}, \field{LcCfRefuse},
\field{LcCfReconcileClaimed}, \field{LcCfAbort} and \field{LcCfExitNoRetry}, plus
\field{LcDeleteSession}. \emph{sweep} is \emph{delete} with \field{LcCfAbort} and
\field{LcCloseUnlocked}; \emph{sweep+ensure} adds \field{LcEnsureCreate},
\field{LcCfRefuse}, \field{LcCfReconcileClaimed}, \field{LcCfExitNoRetry} and
\field{LcCloseLocked}. \emph{claim+delete+close} keeps every operation
except the environment operations, \field{LcEnsureRefresh},
\field{LcEnsureLineageOnly}, \field{LcCreateCli} and \field{LcPurgeSession}.

\paragraph{Results.} Sessions, missions, delegations are the cardinalities
of the three given sets. ``Switches'' lists those set to \field{ztrue}.
The last column is the number of states ProB analysed.

{\scriptsize
\begin{center}
\begin{tabularx}{\textwidth}{@{}l l l X X l r@{}}
\toprule
\textbf{Run} & \textbf{Check} & \textbf{S/M/D} & \textbf{Switches} & \textbf{Outcome} & \textbf{Slice} & \textbf{States} \\
\midrule
\texttt{L\_lin\_w} & LineageStable & 3/2/2 & none & counterexample & lineage & 137 \\
\texttt{L\_lin\_w\_nogone} & LineageStable & 3/2/2 & none; \field{--parent-gone} disabled & counterexample & lineage & 599 \\
\texttt{L\_lin\_order} & LineageStable & 3/2/2 & DeleteOrder & holds, all states visited & lineage & 21953 \\
\texttt{L\_lin\_cycle} & LineageStable & 3/2/2 & CycleCheck & counterexample & lineage & 95 \\
\texttt{L\_lin\_all} & LineageStable & 3/2/2 & DeleteOrder, CycleCheck & holds, all states visited & lineage & 5489 \\
\texttt{L\_acyc\_w} & Acyclic & 3/2/2 & none & counterexample & lineage & 2 \\
\texttt{L\_acyc\_w\_nogone} & Acyclic & 3/2/2 & none; \field{--parent-gone} disabled & counterexample & lineage & 4856 \\
\texttt{L\_acyc\_order} & Acyclic & 3/2/2 & DeleteOrder & counterexample & lineage & 2 \\
\texttt{L\_acyc\_cycle} & Acyclic & 3/2/2 & CycleCheck & holds, all states visited & lineage & 15029 \\
\texttt{L\_acyc\_all} & Acyclic & 3/2/2 & DeleteOrder, CycleCheck & holds, all states visited & lineage & 5489 \\
\texttt{L\_wf\_w} & WellFormed & 3/2/2 & none & holds, all states visited & lineage & 50654 \\
\texttt{L\_wf\_all} & WellFormed & 3/2/2 & DeleteOrder, CycleCheck & holds, all states visited & lineage & 5489 \\
\texttt{P\_w} & NoLiveSessionPurged & 3/2/2 & none & counterexample & purge & 7817 \\
\texttt{P\_w\_nocorrupt} & NoLiveSessionPurged & 3/2/2 & none; no \field{LcCorruptRoster} & counterexample & purge & 8633 \\
\texttt{P\_stalelive} & NoLiveSessionPurged & 3/2/2 & StaleLive & counterexample & purge & 7817 \\
\texttt{P\_corrupt} & NoLiveSessionPurged & 3/2/2 & CorruptPurge & counterexample & purge & 8957 \\
\texttt{P1\_all} & NoLiveSessionPurged & 1/2/2 & StaleLive, CorruptPurge & holds, all states visited & purge & 817 \\
\texttt{P1\_wf} & WellFormed & 1/2/2 & StaleLive, CorruptPurge & holds, all states visited & purge & 817 \\
\texttt{P2\_all} & NoLiveSessionPurged & 2/2/2 & StaleLive, CorruptPurge & no counterexample in states explored; not exhausted & purge & 40000 \\
\texttt{P\_all} & NoLiveSessionPurged & 3/2/2 & StaleLive, CorruptPurge & no counterexample in states explored; not exhausted & purge & 40000 \\
\texttt{Fa\_w} & ReconcileOnExit & 2/2/2 & none & counterexample & claim & 609 \\
\texttt{Fa\_refuse} & ReconcileOnExit & 2/2/2 & RefuseReconciles & counterexample & claim & 2758 \\
\texttt{Fa\_all} & ReconcileOnExit & 2/2/2 & RefuseReconciles, InertBinding & holds, all states visited & claim & 24481 \\
\texttt{Fa\_orph\_w} & NoOrphanOpenSkeleton & 2/2/2 & none & counterexample & claim & 14212 \\
\texttt{Fa\_orph\_nosup} & NoOrphanOpenSkeleton & 2/2/2 & all but Supersede & counterexample & claim & 9640 \\
\texttt{Fa\_orph\_all} & NoOrphanOpenSkeleton & 2/2/2 & all nine & holds, all states visited & claim & 16681 \\
\texttt{Fa\_wf\_all} & WellFormed & 2/2/2 & all nine & holds, all states visited & claim & 16681 \\
\texttt{Fb2\_w} & NoOrphanOpenSkeleton & 2/2/2 & none & counterexample & delete & 248 \\
\texttt{Fb2\_del} & NoOrphanOpenSkeleton & 2/2/2 & DeleteCloses & holds, all states visited & delete & 2561 \\
\texttt{Fc2\_w} & SweepComplete & 2/2/2 & none & counterexample & sweep & 1210 \\
\texttt{Fc2\_post} & SweepComplete & 2/2/2 & PostCheck & counterexample & sweep & 1195 \\
\texttt{Fc2\_del} & SweepComplete & 2/2/2 & DeleteCloses & counterexample & sweep & 1844 \\
\texttt{Fc2\_both} & SweepComplete & 2/2/2 & DeleteCloses, PostCheck & holds, all states visited & sweep & 9569 \\
\texttt{Fc\_w} & SweepComplete & 2/2/2 & none & counterexample & sweep+ensure & 14163 \\
\texttt{Fc\_del} & SweepComplete & 2/2/2 & DeleteCloses & counterexample & sweep+ensure & 24004 \\
\texttt{F2\_orph\_all} & NoOrphanOpenSkeleton & 2/2/2 & all nine & no counterexample in states explored; not exhausted & claim+delete+close & 300000 \\
\texttt{F2\_leak\_all} & ReconcileOnExit & 2/2/2 & all nine & no counterexample in states explored; not exhausted & claim+delete+close & 300000 \\
\texttt{F2\_sweep\_all} & SweepComplete & 2/2/2 & all nine & no counterexample in states explored; not exhausted & claim+delete+close & 300000 \\
\texttt{F3\_leak\_all} & ReconcileOnExit & 3/2/2 & all nine & no counterexample in states explored; not exhausted & claim+delete+close & 60000 \\
\texttt{F3\_orph\_all} & NoOrphanOpenSkeleton & 3/2/2 & all nine & no counterexample in states explored; not exhausted & claim+delete+close & 60000 \\
\texttt{F3\_sweep\_all} & SweepComplete & 3/2/2 & all nine & no counterexample in states explored; not exhausted & claim+delete+close & 60000 \\
\texttt{F3\_wf\_all} & WellFormed & 3/2/2 & all nine & no counterexample in states explored; not exhausted & claim+delete+close & 60000 \\
\bottomrule
\end{tabularx}
\end{center}
}




\paragraph{Verdicts.}
\begin{itemize}
  \item \textbf{LineageStable} holds. With \field{fixDeleteOrder}, at three
    sessions, all 21\,953 states of the \emph{lineage} slice are visited with no
    counterexample (\texttt{L\_lin\_order}); with both lineage amendments,
    5\,489 states (\texttt{L\_lin\_all}). Without it, the shortest
    counterexample is C2. \field{fixCycleCheck} is not needed for it.
  \item \textbf{Acyclic}, in the strong form, holds with \field{fixCycleCheck}
    alone: 15\,029 states visited (\texttt{L\_acyc\_cycle}), and 5\,489 with both
    lineage amendments. Without it the counterexample is C3: one step with
    \field{--parent-gone} (a session its own parent) and six without it. The walk of \S2.3.1 may
    therefore keep its visited set as a defence only.
  \item \textbf{NoLiveSessionPurged} holds with \field{fixStaleLive} and
    \field{fixCorruptPurge}, exhaustively at one session (817 states) and bounded at two and three sessions (see the table).
    Each amendment alone leaves the other's counterexample
    (\texttt{P\_stalelive}, \texttt{P\_corrupt}).
  \item \textbf{NoOrphanOpenSkeleton} holds in its literal form with
    \field{fixDeleteCloses} and \field{fixSupersede}: 16\,681 states of the
    \emph{claim} slice and 2\,561 of the \emph{delete} slice, exhaustively; and
    in the full two-session machine and the three-session machine no
    counterexample in the first 300\,000 and 60\,000 states, which is not exhaustive.
    With \field{fixRefuseReconciles} and \field{fixInertBinding} the exit
    postcondition \field{ReconcileOnExit} holds as well (24\,481 states, all
    visited). \field{SweepComplete} holds with \field{fixDeleteCloses} and
    \field{fixPostCheck} (9\,569 states of the \emph{sweep} slice, all visited).
  \item \textbf{No invariant is left with an open counterexample.} The
    limit is coverage, not a counterexample: the full \field{claim-forked}
    machine is exhausted only slice by slice.
\end{itemize}


\subsection{Consequences for the design document}
\label{sec:lc-consequences}

With every switch on, the operations of this section are the authority.
DES-078 is to be brought into line at these places.

\begin{itemize}
  \item \S2.1.1, \field{ensure} rule~1 and \S2.3.1: add the cycle refusal of
    C3. The sentence ``a cycle cannot be written through \field{ensure}
    except through a \field{--parent-gone} parent'' is false and goes.
  \item \S2.1.1 and \S2.6, \S5 point~1, \field{deleteFiles}: the order is
    close named skeletons, active-mission claim, delegation binding,
    pending record, last-seen, roster, then \fpath{parent-session} last
    (C2, C5). \field{Repair} is unchanged and now covers a failure at the
    last step.
  \item \S5 point~2 item~2 and item~5: a refusal reconciles (C1); the claimed
    case closes the skeleton unless the binding is live, which replaces the
    round-7 paragraph on ``the record and the binding name the same
    delegation'' (C4); the mission is re-read once the skeleton is on disk
    (C6).
  \item \S2.1.1, \field{isStale} and its test table: row~2 becomes not
    stale (C7); an unreadable roster with a fresh last-seen is not purged
    (C8); a test case for each.
\end{itemize}

\subsection{What this model does not establish}
\label{sec:lc-limits}

Each operation is one atomic step, as in the rest of this document, except
that a \field{claim-forked} write phase and a \field{deleteFiles} are
broken at the points where a process can die or a removal can fail. The
locks are therefore assumed to do what the Non-goals section of \texttt{spec-opencode-integration.tex} says; two
exceptions are modelled, the shared mission lock that excludes a locked
sweep and the absence of any exclusion under the unlocked fallback. The
read phase and the start of the write phase are one step: the window in
which the witness releases its claim between them is accepted by the
design and not examined. A read-phase refusal and the unlocked fallback
are always enabled, which over-approximates the real conditions. The claim
\field{claim-forked} writes for the child is given claim origin; the design
does not say. Time is three values with a threshold of one, and the
participant kinds are three; larger values add no behaviour to
\field{isStale}, which compares only. The instance has three sessions for
the lineage and purge slices and two for the \field{claim-forked} slices
(\S\ref{sec:lc-verification}): the claim slices were also run at three,
bounded, and those runs are marked by their state count.

\subsection{Findings about the toolchain}
\label{sec:lc-findings}

\begin{itemize}
  \item The predicate part of a state schema becomes a guard of every
    operation in \texttt{probcli}'s translation (\S\ref{sec:lc-method}). The
    ``zero invariant violations'' of \S\ref{sec:validation} for
    \field{Registry} are claims about a machine in which a violating
    transition cannot occur; they say nothing about whether one is
    possible. This does not affect the \field{Registry} operations'
    guards, which are the stated protocol, but it does affect the word
    ``invariant''. The state predicates of \field{Registry} should be
    re-checked as goals.
  \item With this section present, \texttt{/z-spec:test} of the whole file
    still explores \field{Registry} (286 states, 3862 transitions at
    setsize~1 before; 286 states after), but promotes every schema of this
    section to an operation. Its transition count rises to 14977 and the
    run is no longer certified complete at \field{MAX\_OPERATIONS} 5000. A
    second file for this model, \texttt{docs/spec-opencode-lifecycle.tex},
    would remove the interference; the mission's write-set did not allow it.
\end{itemize}



\end{document}
